\section{Conservativity}
\label{sec:conservativity}

The interpretation of the language in the higher-order model is \emph{conservative} over the first-order
model: every term whose context and type are first-order denotes, up to the isomorphisms relating the two
interpretations, a morphism of the first-order category. The proof is by glueing: the first-order
category embeds into a category of Grothendieck logical relations \citep{fiore-simpson99} over the
higher-order one; the language is interpreted there; and at first-order types the interpretation
collapses back to the first-order category.

\subsection{Predicate systems}
\label{sec:predicate-system}

\subsubsection{Predicate system over $\cat{C}$}

Suppose $\cat{C}$ a category with finite products. A \emph{predicate system} over $\cat{C}$ consists of the
following data:
\begin{itemize}
\item For every object $X$ of $\cat{C}$, a poset $(\Pred(X), \sqsubseteq)$ of predicates on $X$.
\item For every $X$, the structure of a Heyting algebra (without $\bot$) on $\Pred(X)$, i.e.~ top element
   $\top$, meets $\meet$, joins $\join$ and residual $\residual$ satisfying the usual laws.
\item For every morphism $f: X \to Y$ of $\cat{C}$, a Galois connection $\directImage{-}{f} \dashv
\reindex{-}{f} : \Pred(Y) \to \Pred(X)$. The right adjoint $\reindex{-}{f}: \Pred(Y) \to \Pred(X)$ computes
the pullback (reindexing) of a predicate along $f$; the left adjoint $\directImage{-}{f}: \Pred(X) \to
\Pred(Y)$ computes the pushforward (existential image) of a predicate along $f$, with the adjointness property
meaning that
\[\directImage{P}{f} \sqsubseteq Q \iff P \sqsubseteq \reindex{Q}{f} \]
\item For every $X$ and $Y$, a right adjoint $\forall_Y: \Pred(X \times Y) \to \Pred(X)$ to the monotone map
$\reindex{-}{\pi_1}$ called \emph{universal quantification}. \roly{Is there a left adjoint $\exists_Y$?}
\item $\reindex{-}{f}$ is colax (and thus $\directImage{-}{f}$ is lax) with respect to $\top$, $\meet$,
$\residual$ and $\forall_Y$.
\item $\reindex{-}{f}$ is lax (and thus $\directImage{-}{f}$ is colax) with respect to $\join$.
\end{itemize}

\subsubsection{Predicate system over $\Set$}

Suppose $X$ a set. Define a \emph{predicate on $X$} to be a family of (proof-irrelevant) propositions for
every $x \in X$. For any function $f: X \to Y$, define $\reindex{P}{f}(x)$ iff $P(f(x))$ and
$\directImage{P}{f}(y)$ iff there exists $x$ such that $f(x) = y$ and $P(x)$. Define $\mathcal{S}$ to be the
predicate system over $\Set$ which assigns to any set $X$ the set $\Pred(X)$ of propositions on $X$ ordered by
implication, with $\top, \meet, \join$ and $\residual$ defined pointwise and universal quantification defined
as $\forall_Y(P)(x)$ iff $P(x,y)$ for all $y \in Y$.

\subsubsection{Predicates on a presheaf}

Suppose $\cat{C}$ a category and $F: \cat{C}^\op \to \Set$ a presheaf. Define a \emph{predicate on $F$} to be
a family of predicates $P_X \in \Pred_\mathcal{S}(F(X))$ for every object $X$ in $\cat{C}$ such that $P_Y
\sqsubseteq \reindex{P_X}{F(f)}$ for every $f: X \to Y$ in $\cat{C}$.

Now define the predicate system $\mathcal{S}^{\cat{C}}$ over $\Func{\cat{C}^\op}{\Set}$ which assigns to every
presheaf $F$ the poset $\Pred(F)$ of predicates on $F$ ordered pointwise, where $\top, \meet, \join,
\residual$ and $\forall$ in $\Pred(F)$ are all given pointwise, and to every natural transformation $\alpha: F
\to G$, families of pullback and pushforward maps where $\reindex{P}{\alpha}_X = \reindex{P_X}{\alpha_X}$ and
$\directImage{P}{\alpha}_X = \directImage{P_X}{\alpha_X}$.

\subsubsection{Closure operator over a predicate system}

A \emph{closure operator} $C$ over a poset $(X, \sqleq)$ is a lax idempotent monad on $X$, with functoriality,
unit and multiplication given by monotonicity, extensivity ($x \sqleq C(x)$) and idempotence ($C(C(x)) \sqleq
C(x)$).

Suppose $\mathcal{S} = (\Pred, \top, \meet, \join, \residual, \forall)$ a predicate system over a category
$\cat{C}$. A closure operator $C$ over $\mathcal{S}$ is a fibred lax idempotent monad $C$ on the fibration
$\Pred: \cat{C}^\op \to \Poset$ where for every object $X$ the fibre map $C_X: \Pred(X) \to \Pred(X)$ is a
closure operator, $C$ is lax monoidal with respect to $\meet$ (i.e.~$C(P) \meet C(Q) \sqleq C(P \meet Q)$),
and $C$ is lax strong, i.e.~$C(P) \meet Q \sqleq C(P \meet Q)$. $C$ being a fibred functor means it commutes
with any reindexing functor $\reindex{-}{f}$, i.e.~$C_X \comp \reindex{-}{f} \iso \reindex{-}{f} \comp C_Y$.

% Haven't figured this out yet..
%
% \subsubsection{Coproduct closure operator}
%
% Suppose $\cat{C}$ a category with stable coproducts (\secref{stable-coproducts}). Define the fibred functor
% $C$ where for every presheaf $F: \cat{C}^\op \to \Set$ the fibre map $C_F: \Pred(F) \to \Pred(F)$ sends every
% family of predicates $P \in \Pred(F)$ to the family ... for every $X$ in $\cat{C}$

\subsection{The category of logical relations}
\label{sec:conservativity:glueing}

Fix a category $\cat{C}$ with terminal object, finite products and stable finite coproducts (the
extensivity of \citet{fiore-simpson99}), in which the first-order fragment of the language is
interpreted; a category $\cat{D}$ with terminal object,
finite products and coproducts, exponentials, and set-indexed coproducts, in which the full language is
interpreted; and a functor $F : \cat{C} \to \cat{D}$ preserving the terminal object, products and
coproducts.

Predicates live over presheaves on $\cat{C}$: let $G : \cat{D} \to \PSh(\cat{C})$ be the functor
sending $X$ to the presheaf $\cat{D}(F(-), X)$. A presheaf predicate on $G(X)$ thus
assigns to each $y \in \cat{C}$ a predicate on the morphisms $F(y) \to X$, compatibly with
precomposition along $\cat{C}$-morphisms; the predicate system of \secref{predicate-system} provides
reindexing and the connectives for such predicates.

\begin{definition}[Definability]
\label{def:conservativity:definable}
For $x \in \cat{C}$, the predicate $\Definable_x$ on $G(F(x))$ holds at $y \in \cat{C}$ of those $f :
F(y) \to F(x)$ that are images of first-order morphisms: those with $f = F(g)$ for some $g : y \to x$.
\end{definition}

Definability is not closed under the reasoning the interpretation requires; in particular, a morphism
defined by case analysis on a coproduct is definable only up to precomposition with the case split.
Predicates are therefore taken \emph{closed} with respect to a closure operator $\mathbf{C}$
(\secref{predicate-system}), generated by the coproduct structure and the stability assumption on
$\cat{C}$'s coproducts.

The category $\GLR(F)$ of \emph{Grothendieck logical relations} has as objects pairs $(X, \Phi)$ of an
object $X \in \cat{D}$ and a closed predicate $\Phi$ on $G(X)$, and as morphisms $(X, \Phi) \to (Y,
\Psi)$ those $f : X \to Y$ with $\Phi \sqsubseteq \Psi[G(f)]$. Theorem~5.1 of the paper, after \citet{fiore-simpson99},
summarises the structure this category carries: $\GLR(F)$ is bicartesian closed
(here moreover with set-indexed coproducts), each piece built on the corresponding structure of
$\cat{D}$ with an appropriate predicate component; the projection $p : \GLR(F) \to \cat{D}$, $(X,
\Phi) \mapsto X$, strictly preserves it; and the first-order category embeds,
\[
  \GF : \cat{C} \to \GLR(F) \qquad \GF(x) = (F(x), \mathbf{C}\,\Definable_x) \qquad \GF(g) = F(g)
\]
with $\GF$ preserving the terminal object, products and coproducts.

\subsection{Definability}
\label{sec:conservativity:definability}

The point of the construction is that $\GF$ is \emph{full}: morphisms between embedded objects are
already first-order. Fullness is asserted as part of Theorem~5.1 of the paper; here we give the proof.

\begin{lemma}[Fullness]
\label{lem:conservativity:definability}
Every morphism $f : \GF(x) \to \GF(y)$ of $\GLR(F)$ is the image of a first-order morphism: $f = F(g)$
for some $g : x \to y$.
\end{lemma}

\begin{proof}
As a morphism of $\GLR$, $f$ maps $\mathbf{C}\,\Definable_x$ into $\mathbf{C}\,\Definable_y$ reindexed along
$G(f)$; concretely, postcomposition with $f$ sends morphisms satisfying $\mathbf{C}\,\Definable_x$ to
morphisms satisfying $\mathbf{C}\,\Definable_y$. The identity $F(\id_x)$ satisfies $\Definable_x$, so $f = f
\circ F(\id_x)$ satisfies $\mathbf{C}\,\Definable_y$. But $\Definable_y$ is itself closed. The stability of
the coproducts of $\cat{C}$ is used: a morphism definable after case analysis on a coproduct is definable
outright. So $f$ satisfies $\Definable_y$, that is, $f = F(g)$ for some $g : x \to y$.
\end{proof}

The lemma applies only to morphisms between embedded objects, so it remains to show that the
interpretation of every first-order type is isomorphic to one. Interpreting the language in $\GLR(F)$ needs, beyond the structure of
\secref{conservativity:glueing}, only $\Poly$-types (\defref{polynomial-types:has-poly}), which we
construct in \secref{conservativity:glr-poly-types}; their preservation by $\GF$ is assumed here and
analysed in \secref{conservativity:gf-poly-types}. First-order types are also
interpreted directly in $\cat{C}$, which has $\Poly$-types whenever it is a $\Fam$ category
(\propref{polynomial-types:fam-has-poly}); write $\sem{\tau}_{\mathit{fo}}$ for this interpretation. The
two interpretations agree up to the embedding: for first-order $\tau$,
\[
  \sem{\tau} \;\cong\; \GF(\sem{\tau}_{\mathit{fo}})
\]
by induction on $\tau$: the unit, product and sum cases follow from $\GF$'s preservation of the
terminal object, products and coproducts; base types agree by construction, the $\GLR(F)$-model of the
signature being the transport of the $\cat{C}$-model along $\GF$; and $\mu$-types follow from
the assumed preservation together with the isomorphism-invariance of the $\mu$-operator
(\secref{polynomial-types:fam}).

For the $\mu$-free language this is Theorem~5.2 of the paper; the version here adds the
$\Poly$-types hypotheses.

\begin{theorem}[Syntactic definability]
\label{thm:conservativity:syntactic-definability}
Assume $\GF$ preserves $\Poly$-types. Then for every term $\Gamma \vdash t : \tau$
with $\Gamma$ and $\tau$ first-order, there is a first-order morphism $g : \sem{\Gamma}_{\mathit{fo}} \to
\sem{\tau}_{\mathit{fo}}$ with $F(g)$ equal to the underlying $\cat{D}$-morphism of $\sem{t}$, modulo the
isomorphisms above.
\end{theorem}

\begin{proof}
Composing $\sem{t} : \sem{\Gamma} \to \sem{\tau}$ with the isomorphisms gives a morphism
$\GF(\sem{\Gamma}_{\mathit{fo}}) \to \GF(\sem{\tau}_{\mathit{fo}})$ of $\GLR(F)$, to which
\lemref{conservativity:definability} applies.
\end{proof}

\subsection{$\Poly$-types in $\GLR(F)$}
\label{sec:conservativity:glr-poly-types}

$\GLR(F)$ has $\Poly$-types. We prove this by transferring the construction of
\secref{polynomial-types:fam}: realise families of $\GLR(F)$-objects as set-indexed coproducts, and
obtain $\Poly$-types for $\GLR(F)$ from those of $\Fam(\GLR(F))$. The construction is formalised
(\texttt{fam-mu-realisation}), for any category with the structure listed in
\thmref{conservativity:transfer}; preservation by $\GF$ is the subject of
\secref{conservativity:gf-poly-types}.

\paragraph{Realisation.}
Let $\cat{E}$ be a category with set-indexed coproducts. The functor
\[
  \Sigma : \Fam(\cat{E}) \to \cat{E} \qquad \Sigma(X, \partial X) = \coprod_{x \in X} \partial X_x
\]
sends a family to the coproduct of its fibres and a morphism to the copairing of its fibre maps (the
counit of the free coproduct completion). A \emph{presentation} of an object $E \in \cat{E}$ is a
family together with an isomorphism $\Sigma(X, \partial X) \cong E$. In a $\Fam$ category every
object is canonically presented, set-indexed coproducts being disjoint unions of index sets; in
$\GLR(F)$ presentations are genuine data, since a closed predicate on $G(X)$ need not be determined
by its restrictions to the fibres of $X$.

\paragraph{$\Poly$-types via realisation.}
A polynomial $P$ over $\cat{E}$ embeds as a polynomial $\hat{P}$ over $\Fam(\cat{E})$, each
$\const$-object replaced by its $\eta$-image, and $\Fam(\cat{E})$ has $\Poly$-types by
\propref{polynomial-types:fam-has-poly}, since $\cat{E}$ has a terminal object and finite
products. The $\Poly$-types for $\cat{E}$ are the realisations: $\mu_{\alpha.P}(\delta) :=
\Sigma\, \mu_{\alpha.\hat{P}}(\eta\,\delta)$. For this to yield
\defref{polynomial-types:has-poly}, $\Sigma$ must preserve what the construction uses:
set-indexed coproducts (immediate), the terminal object, and finite products, the last requiring
products to distribute over set-indexed coproducts (a consequence of the exponentials).

\paragraph{Operations and laws.}
These do not simply transport along the resulting isomorphisms: the algebra map of
$\mu_{\alpha.\hat{P}}$ lives at an environment containing the $\Fam(\cat{E})$-object
$\mu_{\alpha.\hat{P}}$ itself, whereas its realisation must live at the presented object
$\Sigma\, \mu_{\alpha.\hat{P}}$. The adjunction $\Sigma \dashv \eta$, where $\eta : \cat{E}
\to \Fam(\cat{E})$ embeds an object as a family over a one-element set, transposes algebras and
catamorphisms between the two categories; what remains is that realisation is invariant under
replacing environment entries by families with isomorphic realisations. We package this as a
property of the polynomial, established by induction.

\begin{definition}[Collapse]
\label{def:conservativity:collapse}
A polynomial $P$ over $\cat{E}$ \emph{collapses} if every family of isomorphisms $\Sigma\,
\hat{\delta}_i \cong \Sigma\, \hat{\delta}'_i$ between the realisations of two environments
induces an isomorphism $\Sigma(\hat{P}(\hat{\delta})) \cong \Sigma(\hat{P}(\hat{\delta}'))$,
compatibly with composition of such families, and naturally with respect to the strong action of
$\hat{P}$.
\end{definition}

\begin{remark}
Compatibility with identities is derivable: naturality at projection families gives an
extensionality principle (collapse depends only on the underlying isomorphisms), extensionality
makes the collapse at identity families idempotent via compatibility with composition, and an
idempotent isomorphism is the identity. Compatibility with composition is not derivable from
naturality: the naturality squares require morphisms of $\Fam(\cat{E})$ along the environment
directions, and none exist along a family of mere isomorphisms of realisations.
\end{remark}

\begin{lemma}
\label{lem:conservativity:collapse}
Every polynomial collapses. For $\mu$-polynomials the induced isomorphism is a morphism of
algebras, and the laws of \defref{conservativity:collapse} follow from uniqueness of
catamorphisms.
\end{lemma}

\begin{theorem}[Transfer]
\label{thm:conservativity:transfer}
Let $\cat{E}$ have set-indexed coproducts, a terminal object, finite products, exponentials and
strong coproducts. Then $\cat{E}$ has $\Poly$-types (\defref{polynomial-types:has-poly}).
\end{theorem}

\begin{proof}[Proof sketch]
Take $\mu_{\alpha.P} := \Sigma\, \mu_{\alpha.\hat{P}}$ at singleton environments. The algebra
map is the realised $\Fam(\cat{E})$ algebra map corrected by collapse
(\lemref{conservativity:collapse}) at the bound entry; the catamorphism transposes the
$\Fam(\cat{E})$ catamorphism along $\Sigma \dashv \eta$. For the $\beta$ and $\eta$ laws, the
interpretations agree across a comparison isomorphism $\Sigma(\hat{P}(\hat{\delta})) \cong
P(\delta)$, and the realised strong action of $\hat{P}$ simulates the strong action that $\cat{E}$
derives from the new structure, naturally in this comparison (induction on $P$; the $\mu$ case is
naturality of collapse). The laws are then the $\Fam(\cat{E})$ laws conjugated by the simulation.
\end{proof}

\begin{corollary}
\label{cor:conservativity:glr-poly}
$\GLR(F)$ has $\Poly$-types, its structure being that of
\secref{conservativity:glueing}, with set-indexed coproducts computed as coproducts of carriers
with joined predicates.
\end{corollary}

\subsection{Preservation of $\Poly$-types by $\GF$}
\label{sec:conservativity:gf-poly-types}

Preservation does not follow from $\GF$'s preservation of the finite structure: initial algebras
map out of themselves, so finite preservation yields only the comparison $\mu_{\alpha.\GF(P)}
\to \GF(\mu_{\alpha.P})$, by folding the $\GF$-image of the algebra. The converse direction
requires $\GF$ to preserve the presentation of $\mu_{\alpha.P}$ as a set-indexed coproduct of its
fibres. We accordingly assume, beyond \secref{conservativity:glueing}, that $\cat{C}$ has stable
set-indexed coproducts and that $F$ preserves them; both hold in the intended instance, $\cat{C}$
being a $\Fam$ category and $F$ acting fibrewise.

Preservation asks only for an isomorphism of objects, so the comparison can be carried out on the
concrete carriers of \propref{polynomial-types:fam-has-poly}, where no initiality machinery is
needed: the tree sorts and fibres of the carrier are computed from the index sets and fibres of the
environment entries and $\const$-objects alone. In particular a $\const$-object and a variable
bound to an equal environment entry contribute identical data. Write $P^{*}$ for the constant-free
skeleton of $P$, with the $\const$-objects $\vec{A}$ of $P$ collected as additional environment
entries, and $\check{\delta}_i$ for the presentation of $\GF(\delta_i)$ induced by the canonical
presentation of $\delta_i$: the family over the index set of $\delta_i$ whose fibres are the
$\GF$-images of its fibres.

\begin{lemma}[TODO]
\label{lem:conservativity:definable-coproducts}
If the set-indexed coproducts of $\cat{C}$ are stable then
$\Definable_{\coprod_i x_i} \sqsubseteq \mathbf{C}\,(\bigvee_i \Definable_{x_i}\langle
\mathsf{in}_i \rangle)$.
\end{lemma}

This is the infinitary analogue of the finite-coproduct closure of
\secref{conservativity:glueing}, and is precisely the extension that would admit built-in list
types. Together with preservation of set-indexed coproducts by $F$ it makes $\GF$ preserve them:
carriers by assumption, predicates by the lemma, the coproducts of $\GLR(F)$ being computed as
coproducts of carriers with joined predicates. In particular $\Sigma\check{\delta}_i \cong
\GF(\delta_i)$.

\begin{lemma}[Skeleton; formalised as \texttt{fam-mu-types-2/skeleton}]
\label{lem:conservativity:skeleton}
In a $\Fam$ category, $\mu_{\alpha.Q}(\delta) \cong \mu_{\alpha.Q^{*}}(\delta, \vec{A})$,
where $\vec{A}$ are the $\const$-objects of $Q$.
\end{lemma}

\begin{proof}[Proof sketch]
A $\const$-object and a variable bound to an equal environment entry give the same carrier data,
so the two carriers agree up to renaming of tree constructors, by tree recursion.
\end{proof}

\begin{lemma}[Carrier comparison, TODO]
\label{lem:conservativity:carrier-comparison}
$\GF(\mu_{\alpha.P^{*}}(\delta, \vec{A})) \cong \Sigma(\mu_{\alpha.P^{*}}(\check{\delta},
\check{A}))$.
\end{lemma}

\begin{proof}[Proof sketch]
The two sides interpret the same constant-free polynomial, and the tree sorts coincide: both are
decorated by the index sets of the environment entries, which the presentations preserve. The
carriers of the fibres over a common sort are finite products of the same $F$-images, identified by
tree recursion using $F$'s preservation of products; the predicate components are compared
sort-wise using \lemref{conservativity:definable-coproducts}, the realisation being computed as a
set-indexed coproduct in $\GLR(F)$.
\end{proof}

\begin{proposition}[TODO]
\label{prop:conservativity:gf-preserves-poly}
$\GF$ preserves $\Poly$-types.
\end{proposition}

\begin{proof}[Proof sketch]
Apply \lemref{conservativity:skeleton} in $\cat{C}$, then
\lemref{conservativity:carrier-comparison}, then collapse (\lemref{conservativity:collapse}) at
the constant-free $P^{*}$, at the isomorphism family given by $\Sigma\check{\delta}_i \cong
\GF(\delta_i) \cong \Sigma(\eta\,\GF\,\delta_i)$ and likewise at the entries for
$\vec{A}$, and finally the realisation of \lemref{conservativity:skeleton} in $\Fam(\GLR(F))$,
backwards:
\[
  \GF(\mu_{\alpha.P})
  \;\cong\; \GF(\mu_{\alpha.P^{*}}(\delta, \vec{A}))
  \;\cong\; \Sigma(\mu_{\alpha.P^{*}}(\check{\delta}, \check{A}))
  \;\cong\; \Sigma(\mu_{\alpha.P^{*}}(\eta\,\GF\,\delta, \eta\,\GF\,\vec{A}))
  \;\cong\; \Sigma(\mu_{\alpha.\hat{P}}(\eta\,\GF\,\delta))
  \;=\; \mu_{\alpha.\GF(P)}.
\]
The choice of $\Poly$-types for $\cat{C}$ is immaterial: any two structures have componentwise
isomorphic $\mu$-objects, by folds both ways and uniqueness of catamorphisms.
\end{proof}
